\documentclass[11pt]{article}
\usepackage[a4paper,margin=24mm]{geometry}
\usepackage{amsmath,amssymb,mathtools,amsthm}
\usepackage{hyperref,xurl}
\hypersetup{hypertexnames=false}
\newcommand{\C}{\mathbb C}
\newcommand{\tr}{\operatorname{tr}}
\newcommand{\diag}{\operatorname{diag}}
\newcommand{\sgn}{\operatorname{sgn}}
\newcommand{\rank}{\operatorname{rank}}
\newcommand{\im}{\operatorname{im}}
\newcommand{\eps}{\varepsilon}
\newcommand{\one}{\mathbf 1}
\title{The full signed orbit and the original two-Gram metric\\
Independent finite derivation}
\author{}
\date{20 September 2026}
\begin{document}
\maketitle

\paragraph{Scope and correction.}
This note proves the finite-orbit formulas in Sections 5--6 of the
received calculation \cite{Incoming}%
% BEGIN HISTORICAL EXTERNAL CITATION
~Public reconstruction: \href{https://github.com/KokunoYumeto/zeta-function-research-reader/blob/689a8f87e9c12b7a99b01b3123438c6751ee608f/workbenches/splitzero-tandem/continuations/20260920-native-secular-frame/ORBIT_METRIC_DERIVATION.tex\#L38}{OM1--OM40}.%
% END HISTORICAL EXTERNAL CITATION
{}, and carries their metric calculation
through the original two receiving maps.  Its group formulas (31)--(38)
are valid with a positive quotient metric.  Its equations (40)--(44)
use a specified equal-Gram metric.  The original target has instead
the two coordinate Grams $(G_Y,G)$, and the original source has
$(G_X,G_U)$, as displayed in SE13, TC9, RW14 and RW19 of \cite{Fable}%
% BEGIN VERIFIED PUBLIC PROOF LINK
~\href{https://github.com/KokunoYumeto/zeta-function-research-reader/blob/5992ad50c0f2a06921f1fcaded7b4e38097c0739/workbenches/splitzero-tandem/continuations/20260920-original-kernel-mixed-determinants/providers/FABLE_TO_ORIGINAL_CONDUCTOR.tex\#L1967}{[proof SE11]}%
% END VERIFIED PUBLIC PROOF LINK
{}.
We prove the exact formulas for those original metrics below; equality
of the two Grams is neither required nor inferred.
All root labels, inverse-square-root phases and powers of the boundary
parameter are retained.

\section{Literal states and the finite group}
Let
\[
 h(r)=Ar^4+r^3+Br^2+Cr+D,\qquad A\ne0,
 \tag{OM1}
\]
have four distinct complex roots $r_1,\ldots,r_4$, in a fixed order.
Put $d_j=h'(r_j)$, and fix one of the two choices
$\xi_j^2d_j=1$ at each root.  The two original inverse states over
root $r_j$ are $p_j^\pm=e_j\pm o_j$, where
\[
e_j=\begin{pmatrix}
0\\-r_j\\3id_j\\(B-17r_j+Ar_j^2)d_j-2Ad_j^2
\end{pmatrix},\quad
o_j=\xi_j\begin{pmatrix}
1\\-id_j\\Ad_j^2+2r_jd_j\\i(7r_j^2d_j-13d_j^2)
\end{pmatrix}.
\tag{OM2}
\]
These are the coefficient vectors from \cite{Incoming}%
% BEGIN HISTORICAL EXTERNAL CITATION
~Public reconstruction: \href{https://github.com/KokunoYumeto/zeta-function-research-reader/blob/689a8f87e9c12b7a99b01b3123438c6751ee608f/workbenches/splitzero-tandem/continuations/20260920-native-secular-frame/ORBIT_METRIC_DERIVATION.tex\#L38}{OM1--OM40}.%
% END HISTORICAL EXTERNAL CITATION
{}, equation (1);
they are used literally throughout.  Set $E=[e_1\ \cdots\ e_4]$,
$O=[o_1\ \cdots\ o_4]$, $S=[E\ O]$ and
$V=\C^4_\alpha\oplus\C^4_\beta$.  For original branch coefficients
$c_j^\pm$ the exact coordinate maps are
\[
\alpha_j=c_j^++c_j^-,\quad \beta_j=c_j^+-c_j^-,
\qquad c_j^\pm=\tfrac12(\alpha_j\pm\beta_j),
\qquad \sum_{j,\pm}c_j^\pm p_j^\pm=S(\alpha,\beta).
\tag{OM3}
\]
In particular $\|\alpha\|_2^2+\|\beta\|_2^2=
2\sum_j(|c_j^+|^2+|c_j^-|^2)$.  Euclidean coefficient norms below
refer to the displayed $(\alpha,\beta)$ coordinates, with this factor
retained when returning to the $c^\pm$ coordinates.

Write $P_\pi e_j=e_{\pi(j)}$ and $D_\eps=\diag(\eps_1,\ldots,\eps_4)$.
The finite group and its representation are
\[
\mathcal G=\{(\eps,\pi):\eps_j\in\{1,-1\},
       \ \prod_j\eps_j=\sgn\pi\},\qquad
R_{\eps,\pi}=\diag(P_\pi,D_\eps P_\pi).
\tag{OM4}
\]
Here the signs $\eps$ label output coordinates; a convention in which
they label input roots is converted by $\eps_{\pi(j)}=\sigma_j$.
This states the exact relation to SM4--6 of \cite{Fable}%
% BEGIN VERIFIED PUBLIC PROOF LINK
~\href{https://github.com/KokunoYumeto/zeta-function-research-reader/blob/5992ad50c0f2a06921f1fcaded7b4e38097c0739/workbenches/splitzero-tandem/continuations/20260920-original-kernel-mixed-determinants/providers/FABLE_TO_ORIGINAL_CONDUCTOR.tex\#L1967}{[proof SE11]}%
% END VERIFIED PUBLIC PROOF LINK
{}.
The product is
$(\eps,\pi)(\delta,\sigma)=(\eps\,(\pi\delta),\pi\sigma)$, with
$(\pi\delta)_j=\delta_{\pi^{-1}(j)}$.  Its sign product is
$\sgn(\pi)\sgn(\sigma)$, so it remains in $\mathcal G$, and direct
matrix multiplication proves that $R$ is a representation.  There
are eight sign vectors for each of the twenty-four permutations:
$|\mathcal G|=192$.

\section{Exact averaging and reconstruction}
Let $G$ be a fixed positive Hermitian Gram on the receiving coefficient
space $\C^4$.  Define
\[
P_0=\tfrac14\one\one^*,\quad P_3=I-P_0,\qquad
\lambda_0=\tfrac14\one^*E^*GE\one,\quad
\lambda_3=\frac{\tr(E^*GE)-\lambda_0}{3},\quad
\lambda_4=\tfrac14\tr(O^*GO).
\tag{OM5}
\]
Then
\[
\frac1{192}\sum_{g\in\mathcal G}R_g^*S^*GS R_g
=D_\lambda:=
\diag(\lambda_0P_0+\lambda_3P_3,\lambda_4I_4).
\tag{OM6}
\]
Here is a direct proof of every coefficient.  For fixed $\pi$, the
admissible signs have prescribed product $\sgn\pi$.  The average of
$\eps_i$ is zero: flipping $i$ and any other coordinate is a
sign-product-preserving involution reversing it.  For $i\ne j$,
the average of $\eps_i\eps_j$ is zero by flipping $i$ and a coordinate
outside $\{i,j\}$.  Thus averaging $D_\eps C D_\eps$ keeps exactly
the diagonal of $C$, and averaging any entry of $B D_\eps$ gives
zero.  The mixed blocks of $S^*GS$ vanish under this sign average.
The odd diagonal entries then average under all permutations to
$\tr(O^*GO)/4$.

For completeness let $A_0=E^*GE$.  The average of $P_\pi^*A_0P_\pi$
has constant diagonal $\tr(A_0)/4$ and constant off-diagonal
$\sum_{i\ne j}(A_0)_{ij}/12$.  It therefore has eigenvalue
$\one^*A_0\one/4$ on $\C\one$ and eigenvalue
$(\tr A_0-\one^*A_0\one/4)/3$ on $\one^\perp$.
This is precisely the even block of OM6.  The argument applies
equally to the even and odd sign cosets; averaging only the
even signs at an odd permutation would not enumerate $\mathcal G$.

Every $\lambda$ in OM5 is strictly positive, including at a singular
odd receiving matrix.  To retain quantitative bounds, a coordinate
covector satisfies
\[
|v_k|^2\le (G^{-1})_{kk}\,v^*Gv.
\tag{OM7}
\]
Indeed this is Cauchy--Schwarz applied to $G^{-1/2}e_k$ and
$G^{1/2}v$.  Since $\sum_jr_j=-1/A$, the second row of $E\one$
is $1/A$.  Further,
$3\lambda_3=\tr(P_3E^*GE P_3)$, whose columnwise second coordinates
are $-(r_j-r_{\rm av})$, where $r_{\rm av}=-1/(4A)$.
The first coordinates of $O$ are $\xi_j$, with
$|\xi_j|^2=1/|d_j|$.  Applying OM7 gives
\[
\lambda_0\ge\frac1{4|A|^2(G^{-1})_{22}},\quad
\lambda_3\ge
\frac{\sum_j|r_j-r_{\rm av}|^2}{3(G^{-1})_{22}},\quad
\lambda_4\ge\frac{\sum_j|d_j|^{-1}}{4(G^{-1})_{11}}.
\tag{OM8}
\]
Distinct roots make the middle numerator positive.  No determinant
of $O$ occurs in this proof.

The orbit observation is the explicitly defined map
\[
\mathcal T:V\longrightarrow(\C^4)^{\oplus192},\qquad
(\mathcal Tv)_g=S R_gv,
\tag{OM9}
\]
where the codomain has Gram $G^{\oplus192}$.  Its adjoint, with the
ordinary coefficient Gram on $V$, is
$\mathcal T^\dagger y=\sum_g R_g^*S^*G y_g$.
OM6 gives $\mathcal T^\dagger\mathcal T=192D_\lambda$ and hence
\[
v=\frac1{192}D_\lambda^{-1}\sum_gR_g^*S^*G(\mathcal Tv)_g.
\tag{OM10}
\]
Define $\mathcal L=(192D_\lambda)^{-1}\mathcal T^\dagger$.
Then
$\mathcal L\mathcal L^\dagger=(192D_\lambda)^{-1}$, so its exact
operator norm is
\[
\|\mathcal L\|_{G^{\oplus192}\to2}
=[192\min(\lambda_0,\lambda_3,\lambda_4)]^{-1/2}.
\tag{OM11}
\]
This proves the incoming noise bound and its sharp constant.
The data in OM9 are all 192 displayed observations.  OM10 is not
an identification of this data space with one observation at one
fixed arithmetic parameter.

\section{The actual quotient observation and its entire kernel}
Retain the original coefficient map $\Psi:\C^4\to W$, the original
observation $\Lambda_W:W\to B$, and its actual positive quotient
Gram $Q$ on $B$.  Put
\[
Z=\Lambda_W\Psi,\quad G_\Lambda=Z^*QZ,\quad
(\mathcal T_\Lambda v)_g=ZS R_gv.
\tag{OM12}
\]
No positive definiteness of $G_\Lambda$ is assumed: it is the
pullback of the original quotient metric.  Substitution in OM5--6
gives exactly
\[
\begin{split}
\lambda_{\Lambda,0}&=\tfrac14\|ZE\one\|_Q^2,\\
\lambda_{\Lambda,3}&=\tfrac13\sum_j\|ZE P_3e_j\|_Q^2,\\
\lambda_{\Lambda,4}&=\tfrac14\sum_j\|ZO e_j\|_Q^2,\\
\mathcal T_\Lambda^\dagger\mathcal T_\Lambda
&=192\diag(\lambda_{\Lambda,0}P_0+
\lambda_{\Lambda,3}P_3,\lambda_{\Lambda,4}I_4).
\end{split}\tag{OM13}
\]
The exact whole kernel therefore is the direct sum of those of the
following three spaces for which the displayed test vanishes:
\[
\begin{array}{c|c|c}
\text{space}&\text{dimension}&\text{vanishing test}\\ \hline
\C\one\oplus0&1&ZE\one=0\\
\one^\perp\oplus0&3&ZE P_3=0\\
0\oplus\C^4&4&ZO=0 .
\end{array}\tag{OM14}
\]
Indeed the squared norm of $\mathcal T_\Lambda v$ is the sum of
three nonnegative terms in OM13.  A term vanishes on its whole
component exactly when its coefficient is zero.  As $Q$ is
positive on the attained quotient, a sum of its squared norms
vanishes exactly when every indicated vector vanishes.
If a different presentation retains a semidefinite $Q$, the tests
must instead be made after $Q^{1/2}$; OM13 itself still holds.

This also gives an explicit reconstruction of all surviving states,
without a missing injectivity assumption.  In $D_{\Lambda}^{+}$
replace each positive $\lambda_{\Lambda,j}$ by its inverse and each
zero one by zero, keeping the same $P_0,P_3,I_4$ blocks.  Then
\[
\frac1{192}D_\Lambda^+
\sum_gR_g^*S^*Z^*Q(\mathcal T_\Lambda v)_g
=\diag(\chi_0P_0+\chi_3P_3,\chi_4I_4)v,
\quad
\chi_j=\begin{cases}1&\lambda_{\Lambda,j}>0,\\0&\lambda_{\Lambda,j}=0.
\end{cases}
\tag{OM15}
\]
This is a direct multiplication using OM13, not an unevaluated
choice of a section.  Its image is the Euclidean orthogonal
complement of the exact kernel OM14.  Its noise norm on the
nonzero image is the reciprocal square root of $192$ times the
smallest positive $\lambda_{\Lambda,j}$.  If all vanish, the
receiver is the zero map, as OM15 states.

\paragraph{An evaluated nonzero observation at the new divisor.}
The regular target $(A,B,C,D)=(1,0,c_*,0)$, $c_*=-60/431$, permits
a concrete test of OM14.  Take the explicitly specified scalar
observation
\[
Z=\ell^T=(208c_*^2,\,5ic_*,\,60,\,-21i),\qquad Q=1.
\tag{OM16}
\]
This is a test observation on the receiving coefficient space;
it is not being substituted for the programme's given $\Lambda_W\Psi$.
With $h=r^4+r^3+c_*r$ and $d=4r^3+3r^2+c_*$, reduction of
\[
208c_*^2+5c_*d+60(d^2+2rd)+21(7r^2d-13d^2)
\]
modulo $h$ is zero.  Thus $\ell^TO=0$ for every fixed choice of
the square-root phases.  Reduction of $\ell^T e(r)/i$ gives
\[
f(r)=\frac{60}{185761}
 (1073621r^3+1680900r^2+490909r-75060).
\tag{OM17}
\]
These two reductions are polynomial identities; they follow by
substituting $r^4=-r^3-c_*r$ until degree at most three remains.
To evaluate squared norms without unproved conjugation assumptions,
all four roots are real.  The cubic $r^3+r^2+c_*$ is negative at
$0$, positive at $-2/3$ (its value is $104/11637$), tends to
$-\infty$ at $-\infty$ and to $+\infty$ at $+\infty$.
Its derivative $r(3r+2)$ is nonzero in each of the three open
intervals between these extrema, so it has three distinct real
roots; the fourth root $0$ of $h$ is distinct.

The sums $p_k=\sum_jr_j^k$ are, for $0\le k\le6$,
\[
\left(4,-1,1,-\frac{251}{431},\frac{191}{431},
             -\frac{131}{431},\frac{41401}{185761}\right).
\tag{OM18}
\]
The first three follow from the elementary symmetric coefficients,
$p_3=-1-3c_*$, and for $k\ge4$ the equation of each root gives
$p_k=-p_{k-1}-c_*p_{k-3}$.  Expanding OM17 and its square using
these values gives
\[
\sum_jf(r_j)=\frac{15870600}{185761},\qquad
\sum_jf(r_j)^2=\frac{348296871672000}{34507149121}.
\tag{OM19}
\]
Since $\ell^Te_j=if(r_j)$, OM13 now evaluates to
\[
\lambda_{\Lambda,0}=\frac{62968986090000}{34507149121},
\qquad
\lambda_{\Lambda,3}=\frac{95109295194000}{34507149121},
\qquad
\lambda_{\Lambda,4}=0.
\tag{OM20}
\]
Consequently this nonzero scalar observation has a complete-orbit
map of rank exactly four: it recovers every even state and kills
the entire odd sector.  The fixed observation has rank one, whereas
the displayed 192-coordinate observation has rank four; OM13--15
prove the exact map relating those statements.

\section{The original two-Gram metric}
Let $H>0$ be the label-coordinate Gram and $G>0$ the
evaluation-coordinate Gram.  The coefficient map and its pullback
are
\[
J=\begin{pmatrix}I&0\\E&O\end{pmatrix},\qquad
M_{H,G}=J^*\diag(H,G)J
=\begin{pmatrix}H+E^*GE&E^*GO\\O^*GE&O^*GO\end{pmatrix}.
\tag{OM21}
\]
The exact original receiving maps are
\[
\mathcal J_W(\alpha,\beta)
 =(Y\alpha,\Psi(E\alpha+O\beta)),\quad
\mathcal J_U(\alpha,\beta)
 =(X\alpha,\Phi(E\alpha+O\beta)).
\tag{OM22}
\]
Thus their original orthogonal direct-sum norms give respectively
\[
(H,G)=(G_Y,G)\quad\hbox{and}\quad(H,G)=(G_X,G_U),
\tag{OM23}
\]
where $G_Y=Y^*\langle\, ,\,\rangle_{\Gamma,W}Y$ and
$G=\Psi^*\langle\, ,\,\rangle_{\Gamma,W}\Psi$, and the source
uses the actual quotient metric in both source copies.  Further,
$T_AX=Y$ and $T_A\Phi=\Psi$ prove the full commuting identity
$(T_A\oplus T_A)\mathcal J_U=\mathcal J_W$ by substitution.
This retains the exact morphism between both original metrics.
No equality of $G_Y$ with $G$, or of $G_X$ with $G_U$, is implied.

In the remainder of this section $O$ is invertible, so that $M_{H,G}$
is positive definite.  This holds on the punctured boundary in
the next section.  At the regular-target odd-frame divisor, the
same matrix is semidefinite with its actual kernel; the positive
comparison problem below is not asserted there.
Put
\[
A_0=H+E^*GE,\quad B_0=O^*GO,\quad C_0=E^*GO,\qquad
\zeta_{H,G}=\|G^{1/2}EH^{-1/2}\|_2 .
\tag{OM24}
\]
Among all positive block-diagonal forms $K=\diag(K_e,K_o)$,
the exact least relative comparison ratio is
\[
\inf_{\substack{K>0\ {\rm block\ diagonal},\ 0<\alpha\le\beta\\
                 \alpha M_{H,G}\preceq K\preceq\beta M_{H,G}}}
\frac{\beta}{\alpha}
=\bigl(\sqrt{1+\zeta_{H,G}^2}+\zeta_{H,G}\bigr)^2.
\tag{OM25}
\]
The infimum is attained.  To prove this, the invertible block
congruence $\diag(A_0^{-1/2},B_0^{-1/2})$ carries $M_{H,G}$ to
$\begin{psmallmatrix}I&T\\T^*&I\end{psmallmatrix}$, where
$T=A_0^{-1/2}C_0B_0^{-1/2}$.  This congruence is used solely
to prove the comparison and is inverted below; it does not
replace either original norm.  The matrix
$U=G^{1/2}OB_0^{-1/2}$ is unitary, so
\[
TT^*=A_0^{-1/2}E^*GEA_0^{-1/2},\qquad
\eta:=\|T\|_2=\frac{\zeta_{H,G}}{\sqrt{1+\zeta_{H,G}^2}}.
\tag{OM26}
\]
For the second equality, the eigenvalues of the first matrix are
the generalized eigenvalues of $(E^*GE,H+E^*GE)$.
An eigenvector of $H^{-1/2}E^*GEH^{-1/2}$ with eigenvalue
$s^2$ gives generalized eigenvalue $s^2/(1+s^2)$, and these
vectors form a basis.  Its largest $s$ is $\zeta_{H,G}$.

Choose unit singular vectors $x,y$ with $x^*Ty=\eta$ real
nonnegative.  The vectors $(x,y)$ and $(x,-y)$ have equal
norm in every block-diagonal form, but their squared norms in
the transformed $M_{H,G}$ are $2+2\eta$ and $2-2\eta$.
Consequently any comparison satisfies
$\beta/\alpha\ge(1+\eta)/(1-\eta)$.
Conversely the identity matrix in these coordinates satisfies
the comparison with
$\alpha=1/(1+\eta)$ and $\beta=1/(1-\eta)$:
the eigenvalues of the transformed $M_{H,G}$ are $1\pm s_j(T)$.
Returning to the original coordinates, this identity matrix is
$K=\diag(A_0,B_0)$.  It attains the lower bound.
Finally $(1+\eta)/(1-\eta)=
(\sqrt{1+\zeta_{H,G}^2}+\zeta_{H,G})^2$, proving OM25.

Finite-group averaging of the full original form is also exact:
\[
\overline M_{H,G}
 :=\frac1{192}\sum_gR_g^*M_{H,G}R_g
 =\diag(\mu_0P_0+\mu_3P_3,\mu_4I_4),
\tag{OM27}
\]
where
\[
\mu_0=\tfrac14\one^*(H+E^*GE)\one,\quad
\mu_3=\frac{\tr(H+E^*GE)-\mu_0}{3},\quad
\mu_4=\tfrac14\tr(O^*GO).
\tag{OM28}
\]
This follows from the entrywise sign and permutation sum already
proved in OM6, applied to the three blocks of OM21.  In particular
\[
\det\overline M_{H,G}=\mu_0\mu_3^3\mu_4^4,\qquad
\det M_{H,G}=|\det O|^2\det H\det G.
\tag{OM29}
\]
The second identity uses $\det J=\det O$ and multiplication of
determinants in the exact congruence OM21.

\section{Original escaping boundary, with all leading constants}
Fix $B,C,D$, assume the cubic $r^3+Br^2+Cr+D$ is squarefree
and $\Delta_0:=3C-B^2\ne0$, and vary $A=t\to0$.
Fix the label and evaluation Grams $H,G$ in OM21; the auxiliary
Fable polynomial varies while those original receiving maps
and their physical metrics are held fixed.
Exactly three roots remain bounded: the implicit-function theorem
at each simple cubic root supplies those analytic branches.
For the fourth root, set $v=1/r$.  Its equation is
$t+v+Bv^2+Cv^3+Dv^4=0$, whose derivative with respect to $v$
at $(t,v)=(0,0)$ is one.  Thus
\[
v=-t-Bt^2+O(t^3),\quad
r_\infty=-t^{-1}+B+O(t),\quad
d_\infty=-t^{-2}+4Bt^{-1}+O(1).
\tag{OM30}
\]
An analytic inverse square root on a small punctured disk is
$\xi_\infty=it(1+2Bt+O(t^2))$; the opposite choice negates
the odd column, and does not change any ensuing Gram constant.
Substitution in every coordinate of OM2 gives
\[
e_\infty=-20t^{-3}e_4+O(|t|^{-2}),\qquad
o_\infty=20t^{-3}e_4+O(|t|^{-2}).
\tag{OM31}
\]
For the fourth mean coordinate, the two leading contributions
are $18t^{-1}(-t^{-2})$ and $-2t(t^{-4})$.
For the fourth odd coordinate, $7r_\infty^2d_\infty-
13d_\infty^2=-20t^{-4}+O(t^{-3})$, multiplied by
$i\xi_\infty=-t+O(t^2)$.  The remaining coordinates have
smaller pole order, and the other three columns are bounded.
This proves the constants and phases in OM31.

Let $j_\infty$ be the fixed label of this column.
The rank-one leading matrix gives
\[
E=-20t^{-3}e_4e_{j_\infty}^T+O(|t|^{-2}),\qquad
\zeta_{H,G}
=20\sqrt{G_{44}(H^{-1})_{j_\infty j_\infty}}\,
 |t|^{-3}(1+O(|t|)).
\tag{OM32}
\]
The operator norm of a rank-one matrix $uv^*$ is
$\|u\|_2\|v\|_2$, which gives the coefficient here.
The norm is Lipschitz in the matrix entries, so the matrix
remainder yields the stated relative remainder.  Since
$(\sqrt{1+z^2}+z)^2=4z^2(1+O(z^{-2}))$ for real $z\to\infty$,
OM25 becomes
\[
\boxed{\inf\frac{\beta}{\alpha}
=1600G_{44}(H^{-1})_{j_\infty j_\infty}
 |t|^{-6}(1+O(|t|)).}
\tag{OM33}
\]

For a self-contained determinant calculation at this boundary,
reduce the four polynomial components of $o_j/\xi_j$ modulo
$h$, using $r^4=-(r^3+Br^2+Cr+D)/t$.
The constant first row leaves the following determinant after
the factors $-i$ and $i$ cancel:
\[
\mathsf R(t)=\begin{pmatrix}
2B&3&4t\\
4tBC-8tD-7C&4tB^2-2tC-16t^2D-5B&4tB-8t^2C-3\\
20C/t+76D-52BC&20B/t+5C-52B^2+208tD&20/t-66B+104tC
\end{pmatrix}.
\tag{OM34}
\]
Every entry is a coefficient in the unchanged power basis
$1,r,r^2,r^3$.  Expansion of its leading term gives
\[
\det\mathsf R(t)
=\frac{20}{t}
\det\begin{pmatrix}2B&3&0\\-7C&-5B&-3\\C&B&1\end{pmatrix}
 +O(1)
=\frac{80(3C-B^2)}t+O(1).
\tag{OM35}
\]
If $V_r=\prod_{j<k}(r_k-r_j)$, then
$\prod_jd_j=t^4V_r^2$.  Hence
$(V_r\prod_j\xi_j)^2=t^{-4}$ and
$V_r\prod_j\xi_j=\eps_{\rm br}t^{-2}$ for a locally fixed
$\eps_{\rm br}\in\{1,-1\}$.  Evaluation of the reduced
coefficient matrix at the four roots therefore proves
\[
\det O=\eps_{\rm br}t^{-2}\det\mathsf R(t),\qquad
|\det O|^2=6400|\Delta_0|^2|t|^{-6}(1+O(|t|)).
\tag{OM36}
\]
As $\Delta_0\ne0$, $O$ is invertible for every sufficiently small
nonzero $t$; thus the positivity assumption used in OM25 has
been established on this boundary.

The same single escaping column in OM31 gives
\[
\mu_0=100G_{44}|t|^{-6}(1+O(|t|)),\quad
\mu_3=100G_{44}|t|^{-6}(1+O(|t|)),\quad
\mu_4=100G_{44}|t|^{-6}(1+O(|t|)).
\tag{OM37}
\]
In detail the leading terms of $\tr E^*GE$ and
$\one^*E^*GE\one$ are both $400G_{44}|t|^{-6}$;
division by four gives the first coefficient, and
$(400-100)/3=100$ gives the second.
The odd trace has the same leading coefficient.  The fixed
entries of $H$ have lower order and remain in OM28.
Multiplication of all eight factors in OM29, followed by OM36,
proves
\[
\boxed{\frac{\det\overline M_{H,G}}{\det M_{H,G}}
=\frac{100^8G_{44}^8}
 {6400|3C-B^2|^2\det H\det G}
 |t|^{-42}(1+O(|t|)).}
\tag{OM38}
\]
The numerator grows with order $48$, and the denominator grows
with order $6$.  This proves the exact exponent $42$ and its
coefficient without imposing group invariance on the old metric.

\paragraph{Back to the two original sides.}
Equations OM23, OM33 and OM38 give the complete target constants
\[
1600G_{44}(G_Y^{-1})_{j_\infty j_\infty},\qquad
\frac{100^8G_{44}^8}
 {6400|3C-B^2|^2\det G_Y\det G},
\tag{OM39}
\]
and the complete source constants
\[
1600(G_U)_{44}(G_X^{-1})_{j_\infty j_\infty},\qquad
\frac{100^8(G_U)_{44}^8}
 {6400|3C-B^2|^2\det G_X\det G_U}.
\tag{OM40}
\]
The corresponding exponents are $-6$ and $-42$ in each case.
The incoming equal-Gram formulas follow exactly by setting
$H=G$ in OM33 and OM38, as the explicitly defined metric in
its equation (40) does.  For the actual programme target and
source, OM39--40 retain the distinct original label Grams.
No estimate of a Hermitian norm in this note assigns a
phase-sensitive arithmetic current sign.

\begin{thebibliography}{9}
\bibitem{Incoming}
User-supplied AI web-session calculation,
\emph{The relevant strand is Erd\H{o}s--Straus/Fable},
received 20 September 2026, equations (1)--(44).
No author name was supplied.
The exact incoming file has SHA-256
\path{57efa66b73bb4ce951b1d17ec4170b0f7bf1a00af3718e56faa71e7693ec2717}.
Its original pinned working-source commit is
\path{ab45e221579a0d48ce2885b4ecdf1a6aefc24a20}.
This note reads its entire text and independently proves the
orbit and metric claims; it corrects the native metric identification.


% BEGIN HISTORICAL EXTERNAL CITATION
Public reconstruction: \href{https://github.com/KokunoYumeto/zeta-function-research-reader/blob/689a8f87e9c12b7a99b01b3123438c6751ee608f/workbenches/splitzero-tandem/continuations/20260920-native-secular-frame/ORBIT_METRIC_DERIVATION.tex\#L38}{OM1--OM40}. These links identify the public derivation, not the private working correspondence cited above.
% END HISTORICAL EXTERNAL CITATION
\bibitem{Fable}
Split-Zero programme,
\emph{Fable to the original conductor: Inverse branches, signed
states, and boundary action},
\href{https://github.com/KokunoYumeto/zeta-function-research-reader/blob/5992ad50c0f2a06921f1fcaded7b4e38097c0739/workbenches/splitzero-tandem/continuations/20260920-original-kernel-mixed-determinants/providers/FABLE_TO_ORIGINAL_CONDUCTOR.tex\#L1}{\texttt{FABLE\_TO\_ORIGINAL\_CONDUCTOR.tex}},
accepted local source snapshot, 20 September 2026.
Exact snapshot SHA-256
\path{9629ec8ff32c069a7329c012687fd8bcfd7d5c9437a94fcfba1d4725edbec764}.
The exact sections read for this derivation are SE11--18,
SM3--7, TC8--9 and RW14--20, together with the bibliography.
Repository:
\url{https://github.com/KokunoYumeto/zeta-function-research-reader}.
The source attributes the original marked polynomial to
The Clankers' Erd\H{o}s--Straus reader and the preceding
Alp\"oge--Fable example to the human sources retained there.
Those original publications are not represented as independently
read in this derivation.  The finite averaging, exact kernel,
comparison optimization and boundary constants used here have
complete algebraic proofs above; no unquoted human theorem is
being imported to justify them.

% BEGIN VERIFIED PUBLIC PROOF LINK
\href{https://github.com/KokunoYumeto/zeta-function-research-reader/blob/5992ad50c0f2a06921f1fcaded7b4e38097c0739/workbenches/splitzero-tandem/continuations/20260920-original-kernel-mixed-determinants/providers/FABLE_TO_ORIGINAL_CONDUCTOR.tex\#L1967}{Source proof SE11}.
% END VERIFIED PUBLIC PROOF LINK
\end{thebibliography}
\end{document}
